\section{Introduction}
\label{sec:intro}

\IEEEPARstart{L}{ow} Earth orbit (LEO) satellite constellations have moved from
proposal to deployment: mega-constellations now place thousands of satellites
at altitudes of a few hundred kilometers, and standards bodies position them as
a backbone of 6G non-terrestrial networks (NTN), delivering low-latency global
connectivity where terrestrial infrastructure is absent~\cite{mahboob,kodheli2021,3gppntn}. Unlike bent-pipe relays, modern constellations forward
traffic in space over inter-satellite links (ISLs), so a packet may traverse
many on-board IP routers before reaching a ground gateway~\cite{handley,bhattacherjee2022}.

This in-orbit forwarding is exactly what makes security hard. ISLs are
broadcast over open free-space channels, and the published, near-deterministic
orbital geometry lets an adversary with a suitable terminal predict when and
where a link will exist, then eavesdrop, inject, or replay
traffic~\cite{yue,wu2024leosec,du2025leosec,tang2026satdef,giuliari}. At the same time, an on-board
router runs on a radiation-hardened processor with a fraction of the compute
and memory of a terrestrial core router, and it cannot be patched or rebooted
on demand~\cite{denby}. The result is a sharp mismatch: the threat
surface grows with every added hop, while the per-hop security budget is tiny.


Conventional IP security was designed for endpoints with ample resources and
stable associations. IPsec ESP and DTLS protect a flow by maintaining a
per-flow security association (SA) holding keys, selectors, counters, and a
replay window, and by encrypting and authenticating every packet with
heavyweight primitives~\cite{rfc4301,rfc6347}. This model assumes endpoints
with sufficient memory and relatively stable associations; on a constellation
carrying $10^5$--$10^6$ concurrent flows, the aggregate SA state and per-packet
crypto exceed the on-board budget, and the frequent ISL handovers caused by
orbital motion force costly SA re-establishment, with measurement and modeling
studies of satellite IPsec reporting significant handover and renegotiation
costs~\cite{yue}.

Protocols closer to the space and routing layers help but do not close the gap.
At the link layer, the CCSDS Space Data Link Security protocol authenticates and
encrypts a single space link~\cite{ccsds}; it is effective hop-by-hop but
provides neither end-to-end protection across an in-orbit IP path nor any use of
the network-layer routing information that LEO IP forwarding already carries.
Segment Routing over IPv6 (SRv6) brings programmable source routing to the
satellite data plane and is attractive for LEO because a ground controller can
pin a path through a predictable
constellation~\cite{rfc8402,rfc8754,rfc8986,rfc9252}; the SRH defines an HMAC TLV to
authenticate the segment list and recent work proposes richer security
SIDs~\cite{rfc8754,srsurvey,abdelsalam2021}, but the HMAC TLV authenticates the whole segment
list with a full SHA-256 digest that is expensive to verify on every hop, and
existing SRv6 security designs still assume per-flow or per-policy state at
verifying nodes~\cite{abdelsalam2021,tanveer2024srv6sec}.

Lightweight cryptography supplies fast building blocks but not the missing
systems design. A large body of work, including the NIST
lightweight-cryptography standardization and software-friendly keyed PRFs such
as SipHash, offers primitives suited to constrained
hardware~\cite{nistlwc,simon,siphash,lv2025leolwc}, yet these do not by themselves answer
where to keep security state or how to verify packets at line rate on an
on-board router. Finally, the 20-bit IPv6 Flow Label is set by the source,
immutable in transit, and today used only for ECMP load balancing and QoS
classification~\cite{rfc6437,rfc6438}; to our knowledge it has not been
developed as a network-layer security authenticator with a defined derivation
and verification rule, leaving its potential as a security primitive
unexploited.


The key insight is that security state should reside in the packet rather than
on the satellite, while the structured, predictable LEO topology makes per-hop
checks trivial. This insight is realized in \textbf{6FLAS} (IPv6 \textbf{F}low \textbf{L}abel
\textbf{A}uthenticated \textbf{S}egments), a lightweight security mechanism
that (i) repurposes the IPv6 Flow Label as a connectionless flow authenticator
derived from a pre-shared domain key, and (ii) binds it to an SRv6 security
segment so that every on-board router verifies a single short keyed MAC over
the routing header, using only a small per-direction replay window instead of a
per-flow SA table. Because the Flow Label is set by the source and is immutable
in transit by design~\cite{rfc6437}, it is a natural carrier for a stable, source-bound flow-identity token. Because LEO routers forward over a small
number of fixed grid directions, the egress check that normally requires a
routing lookup collapses to a constant-time comparison.


This paper makes four contributions, each addressing a gap above.

\begin{itemize}

\item \textbf{Packet-carried stateless security model for LEO.}
We formalize LEO security on a resource-constrained directed graph and show how
moving security state from on-board storage into the packet meets the
data-origin authentication, integrity, anti-replay, and access-control
requirements without per-flow SAs (\cref{sec:method}).

\item \textbf{Flow-Label-based stateless flow authenticator.}
We define a flow as a five-tuple plus a security-domain identifier, derive the
20-bit Flow Label from a domain key, and give a stateless verification rule that
needs no per-flow table walk (\cref{sec:fl-auth}).

\item \textbf{SRv6 Security SID integration.}
We integrate the authenticator with SRv6 through a Security SID
(Locator$\,|\,$Function$\,|\,$Arguments) that encodes a path security policy and
specializes per-hop checking to the six fixed LEO forwarding directions, while
interoperating with nodes that implement the base \texttt{End} behavior (\cref{sec:srv6-bind}).

\item \textbf{Security argument and reproducible evaluation.}
We analyze resilience to eavesdropping, forgery, man-in-the-middle, replay, and
stateless DoS, and quantify cost with an explicit processing model and an
event-driven latency model of a 1584-satellite constellation
(\cref{sec:security,sec:exp}).

\end{itemize}


On a 1\,GHz space-grade core, 6FLAS verifies a packet in about $306$\,ns per hop,
roughly $48\times$ faster than IPsec ESP and $5.9\times$ faster than an SRv6
HMAC baseline (analytically projected on a 1\,GHz space-grade core); it adds only $8$ bytes per packet ($79\%$ less than IPsec ESP)
and carries zero per-flow on-board state ($100\%$ reduction versus IPsec~ESP; total fixed state is $192$\,B, direction-scoped). Under a $10^6$-flow load its single-core forwarding throughput is about
$48\times$ that of IPsec ESP, and its end-to-end latency stays within the
propagation-dominated regime of the constellation.

The remainder of this paper is organized as follows.
\cref{sec:prelim} gives the LEO,
Flow Label, and SRv6 background, the threat model, and the problem
formulation. \cref{sec:method} presents 6FLAS. \cref{sec:security} analyzes its
security, and \cref{sec:exp} reports the evaluation. \cref{sec:conclusion}
concludes.
